% ============================================================
% Nature Article LaTeX Template — Correct Structure
% Based on: https://www.nature.com/nature/for-authors/formatting-guide
% ============================================================

\documentclass[10pt]{article}

% --- Required packages ---
\usepackage[margin=1in]{geometry}
\usepackage{amsmath,amssymb,booktabs,graphicx,hyperref,xcolor}
\usepackage[numbers,super]{natbib}  % Nature uses superscript numbered refs
\usepackage{setspace}
\doublespacing                       % REQUIRED: Nature demands double-spacing
\usepackage{lineno}
\linenumbers                         % REQUIRED: line numbers for review

% ============================================================
\title{\textbf{The majority bias: standard single-cell workflows systematically eliminate rare subpopulations during batch integration}}

\author{[Author list]}
\date{}

\begin{document}
\maketitle

% ============================================================
% ABSTRACT — Nature limit: 150 words MAX, single paragraph, no subheadings
% ============================================================
\begin{abstract}
% [MUST BE ≤150 WORDS — current draft is ~200+, needs compression]
Single-cell batch integration harbors a systematic ``majority bias''---implicit assumptions across nine pipeline steps that collectively suppress rare subpopulation signals to $\sim$0.06\% retention. We show that HVG selection excludes 60--100\% of rare-type markers, PCA compresses rare signals by $>$99\%, and Harmony overcorrects by $\sim$5.6-fold. We introduce RASI, improving neighborhood preservation from 0.577 to 0.918 (immune atlas, 2.26M cells) and 0.377 to 0.903 (pancreas), while improving batch mixing. Subpopulation destruction is scale-dependent: pDC are safe at 5 batches but severely disrupted at 103. Wet-lab validation confirms a TIGIT$^+$CCR8$^-$ Treg community disrupted by integration exhibits potent immunosuppression (CD69 4.1$\times$$\downarrow$, IFN-$\gamma$ 2.7$\times$$\downarrow$, killing 3.3$\times$$\downarrow$). Protecting rare subpopulations requires rethinking the entire workflow, not just the integration algorithm.
% [~120 words — within limit]
\end{abstract}

% ============================================================
% INTRODUCTION — Nature: NO section heading, just start writing
% ============================================================
% [Do NOT use \section{Introduction} — Nature starts directly]

Single-cell RNA sequencing has transformed...

% ============================================================
% RESULTS — with subheadings
% ============================================================
\section{Results}

\subsection{The majority bias: a cascade of implicit assumptions}
...

\subsection{RASI: Rare-Aware Single-cell Integration}
...

\subsection{RASI validation across multiple datasets}
...

\subsection{Scale-dependent vulnerability of rare subpopulations}
...

\subsection{Biological validation: TIGIT$^+$CCR8$^-$ Treg functional community}
...

% ============================================================
% DISCUSSION
% ============================================================
\section{Discussion}
...

% ============================================================
% REFERENCES — Nature: BEFORE Methods
% ============================================================
\bibliographystyle{naturemag}
\bibliography{references}

% ============================================================
% METHODS — Nature: AFTER References (this is unique to Nature!)
% ============================================================
\section{Methods}

\subsection{Data}
...

\subsection{RASI pipeline}
...

\subsection{Wet-lab validation}
...

\subsection{Statistical analysis}
...

\subsection{Computational efficiency}
...

% ============================================================
% FIGURE LEGENDS — After Methods
% ============================================================
\section*{Figure Legends}
% [Insert figure legends here — see figure_legends_draft.tex]

% ============================================================
% REQUIRED SECTIONS (Nature mandates all of these)
% ============================================================
\section*{Acknowledgements}
...

\section*{Author contributions}
% Example: X.Y. conceived the study. A.B. performed computational analyses.
% C.D. performed wet-lab experiments. All authors reviewed the manuscript.

\section*{Competing interests}
The authors declare no competing interests.

\section*{Data availability}
Pan-cancer atlas: Zenodo DOI:10.5281/zenodo.10651059.

\section*{Code availability}
RASI source code: [GitHub URL].

% ============================================================
% EXTENDED DATA (up to 10 items, peer-reviewed)
% ============================================================
% Extended Data Figures/Tables go here

\end{document}
